%=============================================================================
\section{Základy odborné práce}
%=============================================================================

\subsection{Vědecká práce a výzkumné metody}

\subsubsection{Základní pojmy}

\begin{itemize}
  \item \textbf{Věda} -- systematická snaha budovat a organizovat znalosti ve formě testovatelných vysvětlení a predikcí o světě
  \item \textbf{Výzkum} -- cílená, systematická a pečlivě prováděná investigace zaměřená na nalezení nových faktů nebo principů
  \item \textbf{Výzkumná metoda} -- soubor technik a procedur používaných k identifikaci a analýze informací
  \item \textbf{Hypotéza} -- testovatelné tvrzení, které může být potvrzeno nebo vyvráceno empirickými daty
  \item \textbf{Teorie} -- ověřené vysvětlení nebo model přijatý vědeckou komunitou
\end{itemize}

\subsubsection{Kvantitativní vs. kvalitativní výzkum}

\begin{center}
\begin{tabular}{p{6cm}p{6cm}}
\toprule
\textbf{Kvantitativní výzkum} & \textbf{Kvalitativní výzkum} \\
\midrule
Numerická data, statistická analýza & Textová/narativní data, interpretace \\
Ověřuje hypotézy & Generuje hypotézy \\
Velký výzkumný vzorek & Malý vzorek, hloubkový pohled \\
Dotazníky, experimenty & Rozhovory, pozorování, případové studie \\
Výsledky lze zobecnit & Výsledky kontextuálně specifické \\
Objektivita & Subjektivita (role výzkumníka) \\
Deduktivní přístup (od teorie k datům) & Induktivní přístup (od dat k teorii) \\
\bottomrule
\end{tabular}
\end{center}

\paragraph{Proč kvantitativní ověřuje a kvalitativní generuje hypotézy?}

Rozdíl plyne z opačného směru uvažování -- \textbf{deduktivního} (od teorie k datům) vs. \textbf{induktivního} (od dat k teorii).

\textbf{Kvantitativní výzkum} vychází z existující teorie nebo dosavadního poznání. Výzkumník nejprve formuluje konkrétní, měřitelnou hypotézu (např. \textit{„Studenti, kteří denně stráví více než 3 hodiny na sociálních sítích, mají nižší průměr známek"}), poté sesbírá numerická data a statisticky otestuje, zda je hypotéza pravdivá. Hypotéza tedy \textbf{existuje před sběrem dat} -- cílem je ji potvrdit nebo vyvrátit.

\textbf{Kvalitativní výzkum} naopak vstupuje do terénu bez pevně dané hypotézy. Výzkumník provádí hloubkové rozhovory, pozoruje nebo analyzuje texty a z dat inductivně \textbf{odhaluje vzorce a témata}. Tato témata pak vedou ke vzniku nových hypotéz nebo teorií. Hypotéza tedy \textbf{vzniká až po sběru dat} jako výstup procesu.

\textbf{Příklad -- výzkum vlivu sociálních sítí na studijní výsledky:}
\begin{itemize}
  \item \textit{Kvalitativní fáze (generuje hypotézu):} Výzkumník provede 20 hloubkových rozhovorů se studenty a zjistí, že mnozí zmiňují „přerušení soustředění notifikacemi" jako hlavní problém. Na základě toho formuluje hypotézu: \textit{„Frekvence notifikací negativně koreluje se studijním výkonem."}
  \item \textit{Kvantitativní fáze (ověřuje hypotézu):} Výzkumník sesbírá data od 500 studentů (počet notifikací/den, GPA) a provede korelační analýzu. Výsledek hypotézu potvrdí nebo vyvrátí.
\end{itemize}

Kvalitativní přístup je tedy typický pro \textbf{explorativní výzkum} -- kdy téma není dobře zmapované a nevíme, co hledat. Kvantitativní přístup nastupuje, když již víme, co testovat, a potřebujeme výsledek zobecnit na větší populaci.

\paragraph{Smíšený výzkum (Mixed Methods)}
Kombinace kvantitativního a kvalitativního přístupu -- využívá výhody obou přístupů.

\subsubsection{Výzkumný proces}

\begin{enumerate}
  \item Identifikace výzkumného problému
  \item Přehled literatury (literature review)
  \item Formulace výzkumných otázek nebo hypotéz
  \item Volba výzkumné metodologie
  \item Sběr dat
  \item Analýza dat
  \item Interpretace výsledků
  \item Prezentace a publikace výsledků
\end{enumerate}

%=============================================================================
\subsection{Struktura odborného textu -- IMRaD}
%=============================================================================

\subsubsection{IMRaD struktura}

\textbf{IMRaD} (Introduction, Methods, Results and Discussion) je standardní struktura vědeckého článku,
zejména v přírodních a sociálních vědách. Odpovídá na čtyři klíčové otázky:

\begin{enumerate}
  \item \textbf{Introduction (Úvod)} -- \textit{Proč byl výzkum proveden?}
    \begin{itemize}
      \item Kontext a motivace výzkumu
      \item Přehled existující literatury (stav poznání)
      \item Identifikace výzkumné mezery (research gap)
      \item Cíl práce a výzkumné otázky/hypotézy
      \item Stručný přehled struktury práce
    \end{itemize}

  \item \textbf{Methods (Metody)} -- \textit{Jak byl výzkum proveden?}
    \begin{itemize}
      \item Výzkumný design a přístup
      \item Popis vzorku a způsob sběru dat
      \item Použité nástroje, techniky a analytické metody
      \item Postup replikovatelný jinými výzkumníky
      \item Etické aspekty výzkumu
    \end{itemize}

  \item \textbf{Results (Výsledky)} -- \textit{Co bylo zjištěno?}
    \begin{itemize}
      \item Objektívní prezentace zjištění bez interpretace
      \item Tabulky, grafy a obrázky s popisky
      \item Statistická data (pokud je výzkum kvantitativní)
      \item Odpovědi na výzkumné otázky
    \end{itemize}

  \item \textbf{Discussion (Diskuse)} -- \textit{Co výsledky znamenají?}
    \begin{itemize}
      \item Interpretace výsledků v kontextu existující literatury
      \item Potvrzení nebo vyvrácení hypotéz
      \item Praktické implikace zjištění
      \item Omezení výzkumu (limitations)
      \item Doporučení pro budoucí výzkum
    \end{itemize}
\end{enumerate}

\paragraph{Doplňující části odborného textu}
\begin{itemize}
  \item \textbf{Abstract (Abstrakt)} -- stručné shrnutí celé práce (150--250 slov); píše se jako poslední
  \item \textbf{Klíčová slova} -- 4--8 termínů pro indexování
  \item \textbf{Závěr (Conclusion)} -- syntéza hlavních zjištění; někdy sloučen s Discussion
  \item \textbf{Literatura (References)} -- seznam použitých zdrojů
  \item \textbf{Přílohy (Appendices)} -- doplňující materiály (dotazníky, rozsáhlá data apod.)
\end{itemize}

\subsubsection{Struktura diplomové práce vs. IMRaD}

IMRaD je formát \textbf{vědeckých a akademických článků} (journal paper, conference paper apod.), zatímco diplomová práce (DP) má rozsáhlejší strukturu odpovídající akademickým požadavkům školy.

\begin{center}
\begin{tabular}{p{5.5cm}p{5.5cm}}
\toprule
\textbf{Diplomová práce (VŠE)} & \textbf{IMRaD (vědecký článek)} \\
\midrule
Titulní strana + prohlášení & -- \\
Poděkování (volitelné) & -- \\
Abstrakt CZ + EN (150--300 slov) & Abstract (100--250 slov) \\
Klíčová slova & Keywords \\
Obsah, seznam obrázků/tabulek & -- \\
\textbf{Úvod} -- cíl, motivace, struktura & \textbf{Introduction} \\
\textbf{Teoretická část} -- přehled literatury (rešerše) & (zahrnut v Introduction) \\
\textbf{Metodika} & \textbf{Methods} \\
\textbf{Analytická/empirická část} & \textbf{Results} \\
\textbf{Výsledky a diskuse} & \textbf{Discussion} \\
\textbf{Závěr} + doporučení & (součást Discussion) \\
Seznam literatury & References \\
Přílohy & Appendices (volitelné) \\
\bottomrule
\end{tabular}
\end{center}

\textbf{Klíčové rozdíly:}
\begin{itemize}
  \item DP obsahuje povinnou \textbf{titulní stranu a prohlášení o samostatném zpracování}
  \item DP má rozsáhlou samostatnou \textbf{teoretickou část} (rešerše literatury); v IMRaD je přehled literatury pouze součástí Úvodu
  \item DP je obvykle \textbf{50--100+ stran}; vědecký článek 8--20 stran
  \item DP musí být schválena vedoucím, IMRaD prochází externím recenzním řízením
\end{itemize}

%=============================================================================
\subsection{Bibliografické citace a odkazování}
%=============================================================================

\subsubsection{Principy citování}

Citování slouží k:
\begin{itemize}
  \item Vzdání uznání autorům citovaných prací
  \item Podpoření vlastních tvrzení autoritativními zdroji
  \item Umožnění čtenáři dohledat zdrojové informace
  \item Prevenci plagiátorství
\end{itemize}

\textbf{Plagiátorství} -- použití cizích myšlenek, textu nebo dat bez řádného odkazu na původního autora.
Zahrnuje: doslovné kopírování, parafrázi bez odkazu, sebeplagiátorství.

\subsubsection{Citační norma ČSN ISO 690}

Mezinárodní norma pro bibliografické citace platná v ČR. Definuje povinné a nepovinné prvky
bibliografického záznamu a způsob jejich uspořádání.

\paragraph{Typy citačních stylů}
\begin{itemize}
  \item \textbf{Harvardský styl (autor-rok)} -- odkaz v textu: (Novák, 2020); seřazení dle autora
  \item \textbf{Číselný styl (Vancouver/IEEE)} -- odkaz v textu: [1]; seřazení dle pořadí výskytu;
    IEEE standard hojně využívaný v informatice a elektrotechnice
  \item \textbf{APA 7 (American Psychological Association)} -- preferovaný styl na VŠE Praha;
    sociální a behaviorální vědy; podrobně níže
  \item \textbf{MLA (Modern Language Association)} -- hojně používaný v humanitních vědách;
    v textu: (Novák 45) -- příjmení + strana bez čárky
  \item \textbf{Chicago} -- dvě varianty: \textit{Author-Date} (podobná APA) nebo \textit{Notes and Bibliography}
    (poznámky pod čarou); hojně v historii a humanitních vědách
  \item \textbf{Oxford (OSCOLA)} -- právní vědy; poznámky pod čarou s čísly
\end{itemize}

\paragraph{APA 7 -- detailní přehled}

APA 7. vydání (2020) je aktuální standard. Klíčové prvky:

\begin{itemize}
  \item \textbf{Odkaz v textu:} (Příjmení, rok) nebo Příjmení (rok) uvádí\ldots
    \begin{itemize}
      \item Dva autoři: (Novák \& Svoboda, 2020)
      \item Tři a více: (Novák et al., 2020) -- od 1. citace
    \end{itemize}
  \item \textbf{Bibliografický záznam -- monografie:}\\
    Příjmení, I. (Rok). \textit{Název knihy}. Nakladatel. \texttt{https://doi.org/xxx}
  \item \textbf{Bibliografický záznam -- článek v časopise:}\\
    Příjmení, I. (Rok). Název článku. \textit{Název časopisu}, \textit{ročník}(číslo), strany. \texttt{https://doi.org/xxx}
  \item \textbf{Novinky APA 7 oproti APA 6:}
    \begin{itemize}
      \item Uvádí se až 20 autorů (dříve 6, pak ``et al.'')
      \item DOI jako hyperlink (\texttt{https://doi.org/...}) povinný, pokud existuje
      \item Místo vydání se \textbf{neuvádí} (dříve povinné)
      \item Přidány pokyny pro citace podcastů, TikTok, Twitter aj.
    \end{itemize}
  \item \textbf{Seznam literatury:} abecedně dle příjmení, visutý odsek (hanging indent) 1,27 cm
\end{itemize}

\paragraph{Struktura bibliografického záznamu (ČSN ISO 690)}

Základní prvky záznamu pro monografii:
\begin{enumerate}
  \item Příjmení a jméno autora
  \item Název díla
  \item Vydání (pokud není první)
  \item Místo vydání
  \item Nakladatel
  \item Rok vydání
  \item ISBN/ISSN
\end{enumerate}

\textbf{Příklady záznamů:}
\begin{itemize}
  \item \textit{Monografie:} NOVÁK, Jan. \textit{Název knihy}. 2. vyd. Praha: Nakladatelství, 2020. ISBN 978-80-000-0000-0.
  \item \textit{Článek v časopise:} NOVÁK, Jan. Název článku. \textit{Název časopisu}. 2020, \textbf{15}(3), 45--67. ISSN 1234-5678.
  \item \textit{WWW zdroj:} NOVÁK, Jan. Název dokumentu [online]. 2020 [cit. 2024-01-15]. Dostupné z: https://www.example.com
\end{itemize}

\subsubsection{Primární a sekundární zdroje}

\begin{itemize}
  \item \textbf{Primární zdroje} -- původní výzkum, originální data; vědecké články (peer-reviewed), disertace, zprávy z výzkumu
  \item \textbf{Sekundární zdroje} -- interpretace nebo analýza primárních zdrojů; přehledové články, učebnice, encyklopedie
  \item \textbf{Terciární zdroje} -- kompilace sekundárních zdrojů; bibliografie, databáze
\end{itemize}

\subsubsection{Recenzní řízení (Peer Review)}

\textbf{Peer review} (recenzní řízení) je proces, při němž nezávislí odborníci ze stejného oboru posuzují vědeckou práci \textbf{před} jejím přijetím k publikaci. Cílem je zajistit kvalitu, vědeckou správnost a relevanci výsledků.

\paragraph{Jak recenzní řízení probíhá -- krok za krokem}
\begin{enumerate}
  \item \textbf{Podání práce} -- autor odešle rukopis redakci časopisu
  \item \textbf{Desk rejection} -- editor rozhodne, zda je práce vhodná pro daný časopis (tematicky i kvalitativně); nevhodné práce jsou okamžitě odmítnuty
  \item \textbf{Výběr recenzentů} -- editor osloví 2--3 odborníky; recenzenti mohou odmítnout
  \item \textbf{Recenze} -- recenzenti obdrží rukopis a do 4--8 týdnů vypracují posudek
  \item \textbf{Rozhodnutí editora}:
    \begin{itemize}
      \item \textit{Accept} -- přijato bez úprav (vzácné u prvního podání)
      \item \textit{Minor revision} -- drobné úpravy, autor reaguje písemně
      \item \textit{Major revision} -- zásadní úpravy, po revizi znovu recenzováno
      \item \textit{Reject} -- odmítnuto; autor může práci přepracovat a podat jinde
    \end{itemize}
  \item \textbf{Revize a odpověď} -- autor odevzdá opravenou verzi + dokument odpovědí na každý komentář recenzenta
  \item \textbf{Finální přijetí a publikace}
\end{enumerate}

\paragraph{Typy peer review}
\begin{itemize}
  \item \textbf{Single-blind} -- recenzenti znají jméno autora, autor nezná recenzenty
  \item \textbf{Double-blind} -- ani recenzenti neznají autora, ani autor recenzenty (nejčastější v sociálních vědách)
  \item \textbf{Open review} -- identity jsou veřejné; transparentnější, ale může odradit recenzenty
  \item \textbf{Post-publication review} -- recenze probíhá po publikaci (např. preprints na arXiv)
\end{itemize}

\paragraph{Konkrétní příklad}
Autor napsal článek o efektivnosti nové metody výuky programování (kvantitativní studie, n=120 studentů). Podá jej do časopisu \textit{Computers \& Education}. Editor osloví dvě recenzenty z oblasti informatiky a pedagogiky. Po 6 týdnech přijde rozhodnutí \textit{Major revision}: recenzent~1 požaduje rozšíření sekce o metodice a doplnění kontrolní skupiny do experimentu; recenzent~2 navrhuje doplnit přehled literatury o 5 klíčových prací. Autor práci přepracuje a napíše detailní odpovědi (``Response to Reviewers''), kde bodově reaguje na každý komentář. Po druhém kole recenzí je práce přijata.

\subsubsection{Etika vědecké práce}

\begin{itemize}
  \item \textbf{Integrita} -- poctivost v hlášení dat a výsledků
  \item \textbf{Objektivita} -- minimalizace předsudků ve výzkumném designu a interpretaci
  \item \textbf{Transparentnost} -- zveřejňování metodologie a dat
  \item \textbf{Ochrana subjektů výzkumu} -- informovaný souhlas, anonymizace dat
  \item \textbf{Konflikt zájmů} -- deklarování potenciálních konfliktů
\end{itemize}

%=============================================================================
\subsection{Informační etika}
%=============================================================================

\subsubsection{Masonův model -- PAPA}

\textbf{Richard Mason (1986)} identifikoval čtyři základní etické problémy informačního věku (model PAPA). Model slouží jako \textbf{checklist pro etické hodnocení IS/IT projektů} -- pomáhá zjistit, kde mohou vznikat etické problémy při návrhu, nasazení nebo používání informačního systému.

\begin{itemize}
  \item \textbf{Privacy (Soukromí)} -- \textit{Jaké informace smíme o lidech sbírat a ukládat?}
    \begin{itemize}
      \item Otázka: Sbíráme data, která opravdu potřebujeme? Informoval jsme uživatele? Mohou ho odmítnout?
      \item Příklady v IS: HR systém ukládající záznamy o zdravotním stavu zaměstnanců;
        e-shop sledující chování uživatele bez souhlasu; Facebook profiling
      \item Relevantní předpisy: GDPR, zákon o ochraně osobních údajů
    \end{itemize}
  \item \textbf{Accuracy (Přesnost)} -- \textit{Kdo odpovídá za správnost informací a jejich opravu?}
    \begin{itemize}
      \item Otázka: Jsou data v systému aktuální a správná? Co se stane, když jsou chybná?
      \item Příklady v IS: chybná kreditní zpráva znemožní získání půjčky; chyba v registru pacientů
        vede ke špatnému léčení; dezinformace na zpravodajském portálu
      \item Kdo nese odpovědnost za opravu? Systém, provozovatel, nebo uživatel?
    \end{itemize}
  \item \textbf{Property (Vlastnictví)} -- \textit{Kdo vlastní data a SW, a jaká jsou práva k jejich použití?}
    \begin{itemize}
      \item Otázka: Patří data zákazníkovi nebo firmě, která je sesbírala? Platím za SW nebo ho mohu kopírovat?
      \item Příklady v IS: uživatelský obsah nahraný na YouTube vlastní YouTube (dle podmínek);
        softwarové licence (proprietární vs. open source); patenty na algoritmy
      \item Klíčové koncepty: autorské právo, licence (GPL, MIT, proprietární), duševní vlastnictví
    \end{itemize}
  \item \textbf{Accessibility (Přístupnost)} -- \textit{Kdo má právo přistupovat k jakým informacím a za jakých podmínek?}
    \begin{itemize}
      \item Otázka: Je přístup k informacím omezený technologicky, ekonomicky, nebo politicky?
      \item Příklady v IS: placené vědecké databáze (Elsevier) vs. open access;
        internetová cenzura; informace dostupné jen majitelům drahých zařízení (digitální propast)
      \item Relevantní: právo na informace, e-government, přístupnost pro ZTP (WCAG)
    \end{itemize}
\end{itemize}

\paragraph{Jak PAPA použít v praxi}
Při návrhu IS (např. personálního systému, mobilní aplikace, e-shopu) projdete každou oblast:
\begin{enumerate}
  \item \textbf{P}: Která osobní data sbíráme? Máme souhlas? Je to minimální nutné množství?
  \item \textbf{A}: Jak zajistíme přesnost dat? Kdo ověřuje a opravuje chyby?
  \item \textbf{P}: Kdo je vlastníkem dat a kódu? Jaká licencování potřebujeme?
  \item \textbf{A}: Kdo může k systému přistupovat? Není přístup nespravedlivě omezený?
\end{enumerate}

\subsubsection{Digitální propast (Digital Divide)}

Nerovnoměrný přístup k ICT technologiím a internetu:
\begin{itemize}
  \item \textbf{Horizontální propast} -- rozdíly ve vlastnictví zařízení a přístupu k internetu
  \item \textbf{Vertikální propast} -- rozdíly v dovednostech a schopnosti využívat ICT efektivně
  \item Dimenze: ekonomická, geografická (venkov vs. město), generační, genderová
\end{itemize}

\subsubsection{Etické problémy informační společnosti}

\paragraph{Mikroetické problémy (úroveň jednotlivce)}
\begin{itemize}
  \item Kyberšikana a obtěžování online
  \item Závisost na technologiích
  \item Narušení soukromí sociálními sítěmi
  \item Neetické využívání AI (deepfakes, manipulace)
\end{itemize}

\paragraph{Makroetické problémy (úroveň společnosti)}
\begin{itemize}
  \item Dohledová společnost (surveillance capitalism) -- sběr a monetizace osobních dat
  \item Algoritmická diskriminace -- automatizovaná rozhodnutí znevýhodňující skupiny
  \item Digitální monopoly (GAFAM -- Google, Amazon, Facebook, Apple, Microsoft)
  \item Dopady automatizace na trh práce
\end{itemize}

%=============================================================================
\subsection{Žánry odborného textu a citační nástroje}
%=============================================================================

\subsubsection{Žánry odborného textu}

\begin{itemize}
  \item \textbf{Abstrakt} -- nejkratší odborný text; shrnutí celé práce (150--250 slov); vždy v češtině i angličtině
  \item \textbf{Anotace} -- několik řádků; autor vysvětlí, co se za názvem tématu skrývá
  \item \textbf{Resumé} -- delší shrnutí práce (celá kapitola nebo část); někdy v jiném jazyce
  \item \textbf{Recenze} -- zhodnocení vědecké práce; subjektivní pohled; zpravidla do 7 stran
  \item \textbf{Rešerše} -- souopis subjektivně nejdůležitějších bodů a myšlenek odborného textu; slouží k rychlé orientaci v tématu
  \item \textbf{Konspekt} -- písemný záznam obsahu a podstatných myšlenek odborného díla; výpisek z práce, často doplněný komentáři autora
  \item \textbf{Teze} -- zhušťený výklad hlavních myšlenek \textit{teprve připravované} publikace; cca 10 stran
  \item \textbf{Odborný esej} -- slohově vybroušené úvahy; cílem je navrhnout řešení odborné otázky; svobodný útvar s vysokými nároky na odbornost
  \item \textbf{Polemika} -- vědecký spor; vyargumentované hájení vlastního názoru vůči jinému textu
  \item \textbf{Referát} -- zpráva o odborném problému a jeho výsledcích; určena k prezentaci
  \item \textbf{Koreferát} -- navazuje na referát; zaměřen pouze na jeden aspekt problému z předchozího referátu
  \item \textbf{Vědecká stať} -- autor koncipuje vlastní nové řešení problému; plnohodnotný odborný žánr
  \item \textbf{Odborný vědecký článek} -- kratší stať představující výsledky k danému problému
  \item \textbf{Příspěvek do sborníku} -- hlavní myšlenka a závěry jsou autorem předneseny na vědecké konferenci
  \item \textbf{Monografie} -- vědecká kniha věnovaná omezenému tématu do hloubky
  \item \textbf{Učebnice} -- odborné texty určené ke vzdělávání čtenářů
  \item \textbf{Seminární / ročníková práce} -- zadávána studentům SŠ i VŠ; nácvik výzkumných a písemných dovedností
  \item \textbf{Závěrečná práce} -- bakalářská, diplomová, rigorózní (po složení magisterské zkoušky)
  \item \textbf{Disertační práce} -- doktorská práce s původním vědeckým přínosem
\end{itemize}

\subsubsection{Citační manažery}

Software pro správu bibliografických záznamů a automatické generování citací:

\begin{itemize}
  \item \textbf{Zotero} -- open source, zdarma; plugin pro prohlížeč; integrace s Word/LibreOffice
  \item \textbf{Citace.com / Citace PRO} -- česky; oblíbený v českém akademickém prostředí
  \item \textbf{EndNote} -- komerční; rozšířený v přírodních a lékařských vědách
  \item \textbf{Mendeley} -- od Elseviera; sociální síť pro vědce, PDF anotace
  \item \textbf{RefWorks} -- cloudový; oblíbený na amerických univerzitách
\end{itemize}

Funkce citačních manažerů: automatický import metadat, generování citací v různých stylech, správa PDF, spolupráce.

\paragraph{APA styl na VŠE}
APA (American Psychological Association) je preferovaný citační styl na VŠE Praha.
Formát: \texttt{Příjmení, Iniciála. (Rok). \textit{Název}. Vydavatel.}
Odkaz v textu: (Novák, 2020) nebo Novák (2020) uvádí\ldots

\begin{profesorbox}[Otázky profesorů k tématu: Základy odborné práce]
\textbf{Pour Jan, doc. Ing., CSc.}
\begin{itemize}
  \item Jaká je obecná struktura odborného textu?
  \item Co je IMRaD a jak se liší od standardní struktury diplomové práce?
  \item Jak rozlišit primární a sekundární zdroje?
\end{itemize}
\textbf{Řepa Václav, prof. Ing., CSc.}
\begin{itemize}
  \item Jaký je rozdíl mezi kvantitativním a kvalitativním výzkumem?
  \item Jak probíhá recenzní řízení vědeckého článku?
\end{itemize}
\textbf{Maryška Miloš, doc. Ing., Ph.D.}
\begin{itemize}
  \item Co je plagiátorství a jak mu předcházet?
  \item Jaké citační normy znáte?
\end{itemize}
\textbf{Chudán David, Ing., Ph.D.}
\begin{itemize}
  \item Jaký je postup při tvorbě odborné práce?
  \item Co tvoří abstrakt odborného článku?
\end{itemize}
\textbf{Černý Jan, Ing., Ph.D.}
\begin{itemize}
  \item Co je etika a informační etika?
  \item Jak PAPA model pomáhá při etickém hodnocení IS?
\end{itemize}
\textbf{Sigmund Tomáš, doc. Ing., Ph.D.}
\begin{itemize}
  \item Jaké jsou zásady kvantitativního výzkumu?
  \item Jaké jsou zásady kvalitativního výzkumu?
  \item Jak rozlišit kvantitativní a kvalitativní přístup k výzkumné otázce?
\end{itemize}
\textbf{Sklenák Vilém, prof. Ing., CSc.}
\begin{itemize}
  \item Jaké znáte typologie informačních zdrojů?
  \item Jaký je rozdíl mezi primárními a sekundárními zdroji?
  \item Jaká jsou kritéria hodnocení věrohodnosti informačního zdroje?
\end{itemize}
\end{profesorbox}
